% Generated draft: requires mathematical and source review.
\documentclass[11pt]{article}
\usepackage[margin=25mm]{geometry}
\usepackage{fontspec}
\setmainfont{Latin Modern Roman}
\usepackage{amsmath,amssymb,mathtools,unicode-math}
\setmathfont{Latin Modern Math}
\usepackage{enumitem,needspace,xurl,hyperref}
\setlength{\parindent}{0pt}
\setlength{\parskip}{5pt}
\setlength{\emergencystretch}{3em}
\newcommand{\N}{\mathbb{N}}
\newcommand{\Z}{\mathbb{Z}}
\newcommand{\Q}{\mathbb{Q}}
\newcommand{\R}{\mathbb{R}}
\newcommand{\C}{\mathbb{C}}
\newcommand{\F}{\mathbb{F}}
\newcommand{\A}{\mathbb{A}}
\newcommand{\PP}{\mathbb{P}}
\newcommand{\E}{\mathbb{E}}
\newcommand{\eps}{\varepsilon}
\newcommand{\dd}{\,\mathrm{d}}
\begin{document}
\section*{UnsolvedMath ID: 30001786}
\textbf{Source problem code:} catalogue entry
\subsection{Definitions and mathematical statement}
Let
\[
D^*=\{z\in\C:0<|z|<1\},\qquad \omega=f(z)\,\dd z,
\]
where \(f:D^*\to\C\) is holomorphic and has an essential singularity at \(0\). Write its Laurent expansion as
\[
f(z)=\sum_{k\in\Z}a_kz^k,\qquad a_k\in\C,
\]
convergent on \(D^*\). The essential-singularity hypothesis means that infinitely many coefficients \(a_k\) with \(k<0\) are nonzero.

The source denotes by \(L(0)\) the linear-approach space at \(0\). Its question is: can the Laurent coefficients \((a_k)_{k\in\Z}\) determine the rotational components of \(L(0)\)?

This formulation retains unresolved source terminology. The supplied record does not define linear approaches, their equivalence as germs, or the relation defining rotational components. It also does not specify whether \(0\) must correspond to a point in the metric completion for the flat length element \(|f(z)|\,|\dd z|\), how approaches of infinite length are treated, or how zeros of \(f\) are handled. No zero-free hypothesis is supplied. Finally, the requested meaning of ``determine'' is unspecified: the record does not state whether it asks for coefficient criteria for component membership, a description of all components, or further geometric information. These conventions and the intended output must be specified before the question becomes a fully precise characterization problem.
\subsection{Short English statement}
For a holomorphic one-form on the punctured unit disk with an essential singularity at the origin, ask how its Laurent coefficients describe the rotational components of the space of linear approaches to the origin. The supplied record leaves the definitions of these spaces and the precise meaning of the requested characterization unresolved.
\subsection{Sources}
\begin{itemize}
\item Supplied local UnsolvedMath registry record, permanent ID \(30001786\), title: \emph{Laurent Coefficients and Rotational Components of Singular One-Forms}; exact source problem code: \texttt{catalogue entry}. The record attributes the question to the 2011 Oberwolfach workshop report and supplies no individual proposer.
\item \emph{Billiards, Flat Surfaces, and Dynamics on Moduli Spaces} (2011), Oberwolfach Reports, pp.~1361--1427, DOI: \texttt{10.4171/owr/2011/25}. Supplied source URL and literature-source link labeled ``original question'': \url{https://doi.org/10.4171/owr/2011/25}.
\item Supplied literature assessment, checked on \texttt{2026-08-21}: the record labels the problem open and reports that no general Laurent-series characterization was found. This is supplied dated curation, not an independently verified literature assessment.
\item The supplied statement status is \texttt{corrected\_verified}; its accompanying verification note says that the mathematical content was checked against the official report and that notation, context, and question form were normalized without changing the source-level claim. That verification is reported from the record and has not been independently performed here.
\end{itemize}
\end{document}
